\section{The complete low arity filtration for a degree-three Tate stage}
\label{sec:136-q3-low-filtration}

\paragraph{Archive and scope.}
The problem is cote 136, physical pp.~117, 119, 121,
footers 116, 118, 120, scans 117, 119, 121.  This calculation concerns
the proposed preservation of degree-three graph compatibility
after a free degree-three Tate stage.  It gives an exhaustive
filtration of its new third cohomology, with explicit cycle
conditions.  It does not prove that every resulting graph residue
vanishes, and does not shrink the main gate by assertion.
The filtration and all its differentials below are proved here;
no outside theorem or preprint is an input.  The final isolated
Tor observation uses the checked common-skeleton degree-two
primitive lemma ARCH-136-SKELETON-TATE-LIFTS and the elementary
free-unit-disk calculation already checked in
ARCH-136-RELATIVE-DISK-FACTORIZATION.

Let \(S=\Gamma_{\mathbb Z}(T)\), and let \(A\) be a
nonnegative weight-graded special differential PD algebra with
\[
                         H^0(A)=S,\qquad H^1(A)=0
                                                               \tag{136.QF1}
\]
in every weight, with the indicated coefficient map.  At one
simplex of an actual free relative stage, write its new pairs as
\[
 \begin{aligned}
 E&=A\langle t_i,c_i:i\in I\rangle,
       \qquad |t_i|=1,\quad |c_i|=2,\\
 \mathrm{wt}(t_i)&=\mathrm{wt}(c_i)=v_i>0,
       \qquad dc_i=z_i,\quad dt_i=e_i-N_i c_i,\\
 de_i&=N_i z_i,\qquad dz_i=0,\quad N_i\geq1.
 \end{aligned}
                                                               \tag{136.QF2}
\]
The odd variables are exterior and the even variables have free
divided powers.  Every expression uses finitely many variables.
No freeness or integer-torsion-freeness of the form groups of
\(A\) is assumed.  In particular (136.QF1) concerns cohomology,
not the whole group \(A^0\).

Give each new variable arity one and the old algebra arity zero.
Then
\[
               d=d_0+\delta,\qquad
 \begin{cases}
 d_0|_A=d_A,\ d_0t_i=-N_i c_i,\ d_0c_i=0,\\
 \delta|_A=0,\ \delta t_i=e_i,\ \delta c_i=z_i.
 \end{cases}                                                \tag{136.QF3}
\]
Here \(d_0\) preserves arity and \(\delta\) lowers it by
one.  They square to zero and anticommute; on \(t_i\) the
mixed identity is exactly \(de_i=N_i z_i\).  The increasing
arity filtration is finite in every fixed form degree.

\subsection{The first page and its actual representatives}

Put
\[
 L=\bigoplus_i(S/N_i)\langle v_i\rangle,\quad
 V=\bigoplus_iH^2(A)[N_i]\langle v_i\rangle,\quad
 T=\bigoplus_{i<j}(S/g_{ij})\langle v_i+v_j\rangle,
 \qquad g_{ij}=\gcd(N_i,N_j).
                                                               \tag{136.QF4}
\]
The order on the index set is used only to write pairs once.
Define \(\psi:L\to H^3(A)\) by
\(\psi(1_i)=[z_i]\).  This is well-defined because
\(N_i z_i=de_i\).

For \([x_i]\in H^2(A)[N_i]\), take
\(dx_i=0\), \(dy_i=N_i x_i\), with \(|y_i|=1\).
Define
\[
 \Phi:V\longrightarrow H^4(A),\qquad
 ([x_i])_i\longmapsto
                   \left[\sum_i(x_i e_i-y_i z_i)\right].
                                                               \tag{136.QF5}
\]
Every summand is closed.  Changing \(y_i\) by a closed
one-form changes it by an exact one-form, by (136.QF1), and
therefore changes the displayed four-form by a boundary.
Replacing \(x_i\) by \(x_i+du_i\) and \(y_i\) by
\(y_i+N_i u_i\) changes that four-form by \(d(u_i e_i)\).
This proves that \(\Phi\) is a well-defined graded
\(S\)-linear map.

For a pair \(i<j\), write
\(g=g_{ij}\), \(\alpha=N_i/g\), \(\beta=N_j/g\),
and define the arity-two cycle for \(d_0\)
\[
                   \tau_{ij}=\beta t_i c_j-\alpha c_i t_j.
                                                               \tag{136.QF6}
\]
It represents the indicated generator of \(T\), since
\(d_0(t_i t_j)=g\tau_{ij}\).  Its arity-one differential is
\[
 \delta\tau_{ij}
       =\beta z_jt_i-\alpha z_it_j
                         +\beta e_ic_j-\alpha e_jc_i.
                                                               \tag{136.QF7}
\]
This is \(d_0\)-closed.  Its class defines a map
\[
 d_1:T\longrightarrow
       H^4(A\otimes_S K^{(1)}),\qquad
 K^{(1)}=\bigoplus_i[S t_i\xrightarrow{-N_i}S c_i],
                                                               \tag{136.QF8}
\]
where the two terms of each bracket have forms one and two.
Changing a coefficient lift modulo \(g\) in (136.QF6)
changes it by a \(d_0\)-boundary, so (136.QF8) is
well-defined.  For clarity each summand of its target has
the exact sequence
\[
 0\longrightarrow \bigl(H^2(A)/N_iH^2(A)\bigr)\langle v_i\rangle
 \longrightarrow H^4(A\otimes K_i)
 \longrightarrow H^3(A)[N_i]\langle v_i\rangle
 \longrightarrow0.
                                                               \tag{136.QF9}
\]
Thus checking only the old three-class coefficients of
(136.QF7) can miss part of \(d_1\) when \(H^2(A)\ne0\).

\paragraph{ARCH-136-Q3-LOW-FILTRATION.}
There is a natural filtration
\(0\subseteq\mathcal F_0\subseteq\mathcal F_1
\subseteq\mathcal F_2=H^3(E)\), induced by new arity, with
\[
 \begin{aligned}
 \mathcal F_0&\cong H^3(A)/\operatorname{im}\psi,\\
 \mathcal F_1/\mathcal F_0&\cong\ker\Phi,\\
 \mathcal F_2/\mathcal F_1&\cong\ker d_2,
 \qquad d_2:\ker d_1\longrightarrow H^4(A)/\operatorname{im}\Phi.
 \end{aligned}                                                \tag{136.QF10}
\]
Here \(d_2\) is explicit.  Represent \(\xi\in\ker d_1\)
by a sum \(\tau=\sum s_{ij}\tau_{ij}\).  Choose an
arity-one three-form \(h\) with
\[
                         d_0h=-\delta\tau.
\]
Then
\[
                            d_2(\xi)=[\delta h]
                                  \pmod{\operatorname{im}\Phi}.
                                                               \tag{136.QF11}
\]
It is well-defined.  An actual closed three-form with that top
component exists if and only if this class vanishes.  The
filtration is an associated-graded description, not a claimed
canonical direct sum decomposition.

\begin{proof}
We compute the finite filtered complex in the relevant degrees.
For a single arity-one block, the mapping-cone exact sequence
and (136.QF1) give
\[
 H^2(A\otimes K_i)=(S/N_i)\langle v_i\rangle,
 \qquad H^3(A\otimes K_i)=H^2(A)[N_i]\langle v_i\rangle.
                                                               \tag{136.QF12}
\]
The latter class has representative \(x_i t_i+y_i c_i\),
whose actual differential is exactly
\(x_i e_i-y_i z_i\).  This proves (136.QF5) as the
first filtered differential.  The first equality has generator
\(c_i\), whose next differential is \(z_i\), giving
\(\psi\).

The same-variable arity-two block is
\[
                  S(t_i c_i)\xrightarrow{-2N_i}
                                          S\gamma_2(c_i)
                                                               \tag{136.QF13}
\]
in forms three and four.  Its tensor with \(A\) has no
third cohomology: that group is \(H^0(A)[2N_i]=0\).
For two different variables the block has basis \(t_i t_j\)
in form two, \(t_i c_j,c_i t_j\) in form three, and
\(c_i c_j\) in form four.  Choose integers \(a,b\) with
\(a\alpha+b\beta=1\).  The form-three basis
\[
             \tau_{ij},\qquad -a t_i c_j-b c_i t_j
\]
is unimodular.  Its differentials split the block into two
integer Moore blocks, with coefficient \(g\), in forms
\((2,3)\) and \((3,4)\).  After tensoring with \(A\),
(136.QF1) shows that their third cohomology is respectively
\(S/g\) and zero.  Their second cohomology is zero.
Every basis change preserves the original weight.

In new arity three, the only possible form-three monomials
are \(t_i t_j t_k\) with distinct indices.  A cycle's
coefficients must first be closed old zeroforms, hence in
\(S\); the independent terms \(c_i t_jt_k\) in its
differential then force \(N_i\) times each coefficient
to be zero.  They are all zero.  Higher arity cannot occur
in form three.  We have therefore proved the complete first-page
table in total form three:
\[
 \begin{array}{c|cccc}
 \text{new arity}&0&1&2&\geq3\\ \hline
 H^3(d_0)&H^3(A)&V&T&0.
 \end{array}                                                \tag{136.QF14}
\]
In total form two the only positive-arity cohomology is \(L\)
in arity one; there is none in arity two or higher.  These
calculations use actual free integer blocks, not flatness of
the old algebra over \(S\).

It follows already that the old third cohomology is killed
only by \(\psi\), and that the arity-one associated quotient
is \(\ker\Phi\).  There are no further incoming differentials:
their possible sources in total form two and arity at least two
are zero.  The top quotient first has to lie in \(\ker d_1\).
For such a representative, anticommutation shows that
\(\delta h\) in (136.QF11) is a closed old four-form.
Changing \(h\) by a \(d_0\)-closed arity-one three-form
changes its class by an element of \(\operatorname{im}\Phi\).
Changing \(h\) by a \(d_0\)-boundary changes it by an
old boundary.  If \(\tau\) is changed by \(d_0b\), replace
\(h\) by \(h+\delta b\); this changes no \(\delta h\)
because \(\delta^2=0\).  These observations prove all
independence assertions in (136.QF11).

If that quotient class vanishes, alter \(h\) by a closed
arity-one three-form to make \(\delta h\) exact, and choose
an old three-form \(a\) with \(da=-\delta h\).  Then
\(\tau+h+a\) is an actual cycle.  Conversely any cycle
with top component \(\tau\) supplies exactly these equations.
No later outgoing differential can occur from arity two, and
no incoming one exists by the same total-form-two calculation.
This proves (136.QF10).  Equivalently it is the elementary
finite filtration argument with differential (136.QF3);
there is no convergence issue, even for infinitely many cells.
\end{proof}

When \(H^2(A)=0\) exactly, \(V=0\) and (136.QF9)
reduces to its old three-class term.  The first survival
condition for \(\sum s_{ij}\tau_{ij}\) is then the rowwise
vanishing of the classes of the \(t_i\)-coefficients in
\(\sum s_{ij}\delta\tau_{ij}\).  These are sums of
old \([z_j]\), not separate conditions on each pair.
Choose old two-forms \(u_i\) killing those coefficients,
and then old one-forms \(r_i\) killing the remaining closed
two-form coefficients; the latter step uses \(H^2(A)=0\).
For \(h=\sum(u_i t_i+r_i c_i)\), the last obstruction is
the actual old four-class
\[
                              \left[\sum_i(u_i e_i-r_i z_i)\right].
                                                               \tag{136.QF15}
\]
It still has to vanish.  Neither the first survival condition
nor (136.QF15) may be discarded in a proposed generation argument.

\subsection{An isolated Tor generator, and the coefficient-square limitation}

\paragraph{ARCH-136-Q3-ISOLATED-TOR.}
Suppose \(H^2(A)=0\).  If the unmultiplied generator
\(\tau_{ij}\) satisfies \(d_1(\tau_{ij})=0\), it lifts
to a genuine Tor cycle of two actual degree-two torsion classes
in \(E\).  In particular its secondary obstruction vanishes.
This statement for an individual generator does not assert the
same for arbitrary coefficient combinations in \(\ker d_1\).

\begin{proof}
Write \(g,\alpha,\beta\) as above.  The hypothesis says
\(\beta[z_j]=0\) and \(\alpha[z_i]=0\).  Choose
\[
 du_i=-\beta z_j,\qquad du_j=\alpha z_i.
\]
The old two-forms \(e_j+g u_i\) and \(g u_j-e_i\)
are closed.  Choose old one-forms with
\[
                dr_i=e_j+g u_i,\qquad dr_j=g u_j-e_i.
\]
Define
\[
 \begin{array}{ll}
 x_1=\beta c_j+u_i,&y_1=r_i-t_j,\\
 x_2=-\alpha c_i+u_j,&y_2=t_i+r_j.
 \end{array}
 \qquad dx_1=dx_2=0,\quad dy_1=gx_1,\quad dy_2=gx_2.
                                                               \tag{136.QF16}
\]
Then
\[
             z=y_2x_1-y_1x_2,
                    \qquad dz=0,\qquad gz=d(y_1y_2)
                                                               \tag{136.QF17}
\]
is the desired actual cycle.  Its arity-two component is exactly
\(\tau_{ij}\).  The lower-arity calculation in (136.QF10)
shows that any other lift differs, modulo boundaries, by an old
third class.  This proves the assertion.
\end{proof}

In the common-skeleton target setting, the displayed Tor Moore
pair has zero graph-difference residue once the two degree-two
pairs (136.QF16), of individual weights \(v_j\) and \(v_i\),
are each in the fixed degree-two target guard of
\nolinkurl{ARCH-136-SKELETON-TATE-LIFTS}.  A bound only on
their summed weight is not enough for that application.  Indeed choose
common-graph primitives \(da_k=x_k\),
\(db_k=g a_k-y_k\), with \(|a_k|=1\), \(|b_k|=0\).
Their images factor through free unit disks in forms \((1,2)\)
and \((0,1)\), with evaluations agreeing on the graph.
That free algebra has \(H^2=H^3=0\), so the particular
closed Moore pair \((y_1y_2,z)\) has a primitive there.
Its target images have common graph restrictions.  This verifies
the asserted residue for the actual displayed pair; it is not
an exhaustive generation theorem for (136.QF10).  An arbitrary
other lift can differ by an old third class, as proved above;
to conclude its residue also vanishes would additionally require
the old third-cohomology graph-compatibility hypothesis.

Here is a precise limitation even before rowwise combinations.
If only \(s\tau_{ij}\), for a homogeneous coefficient
\(s\in S\), satisfies \(d_1=0\), the analogous construction
starts with
\(du_i=-s\beta z_j\), \(du_j=s\alpha z_i\).
It gives
\[
 x_1=s\beta c_j+u_i,\quad y_1=r_i-s t_j,\qquad
 x_2=-s\alpha c_i+u_j,\quad y_2=s t_i+r_j,
 \quad dy_k=gx_k.
                                                               \tag{136.QF18}
\]
Its Tor cycle has leading term \(s^2\tau_{ij}\), not
\(s\tau_{ij}\).  Multiplication by a positive coefficient
need not be cancellable in \(S/g\).  Thus this argument
does not prove that the original coefficient-multiplied class
has zero residue or zero secondary obstruction.  Sums can also
survive through rowwise cancellations among the \([z_j]\),
without individual pairwise survival.  These are actual remaining
algebraic conditions, not resolved by the isolated construction.

\paragraph{Low repair and the remaining issue.}
A closed one-form of \(E\) is \(a+\sum f_i t_i\).
Its independent new coefficients give \(df_i=0\) and
\(N_i f_i=0\), so \(f_i\in S\) and \(f_i=0\).
Consequently \(H^0(E)=S\) and \(H^1(E)=0\) exactly
under (136.QF1).  For an actual simplicial stage with these
all-weight hypotheses and unchanged coefficient simplex zero,
the supplied RH1 construction has no low cohomology to repair;
its normalized replacement identifies with the existing zeroforms.
Thus the nonlinear third-cohomology issue here is already present
before, and is not changed by, that immediate RH1 step.

To prove general graph-compatibility preservation one must still
control the graph residues of the mixed \(\ker\Phi\) piece
and every surviving coefficient-dependent class of \(\ker d_2\).
If \(H^2(A)\ne0\), it cannot be omitted; a degree-three
repair itself can create new second cohomology.  No preservation
theorem, admissible common source, or strictly smaller main gate
is claimed in this module.

\paragraph{Verification status.}
The author completed a whole saved-proof reading on 22 November
2026, including the final display and scope corrections.  The
parent and reviewer1348 each completed a whole saved-proof reading
and passed both statements at the stated scope.  The parent also
read the complete common-skeleton primitive dependency.  Their
first readings identified the display typo corrected here; no
mathematical conclusion was changed.  The filtration and the
coefficient-square limitation were communicated to the parent
and reviewer123 before saving; those earlier communications are
not counted as whole-proof reviews.  Actual events are recorded
in the check register, with model effort not exposed.
